% =====================================================================
%  8. Marco conceptual
% =====================================================================
\section{Marco conceptual}
\label{sec:marco-conceptual}

Las definiciones de esta sección son \textbf{operacionales}: fijan el sentido en que
cada término se usa en el resto del documento y, cuando corresponde, cómo se mide.
Un término que no pueda definirse de forma que permita decidir si algo cae o no bajo
él, no se usa.

\subsection{Términos del dominio}
\label{sec:conc-dominio}

\begin{xltabular}{\textwidth}{@{}L{3.6cm} X@{}}
\toprule
\textbf{Término} & \textbf{Definición operacional} \\
\midrule
\endfirsthead
\toprule
\textbf{Término} & \textbf{Definición operacional} \\
\midrule
\endhead
\bottomrule
\endfoot

\textbf{Corpus de curso} &
Conjunto cerrado y conocido de documentos asociados a una asignatura, cargados por
el docente responsable. Es la referencia contra la cual se define la corrección. \\

\textbf{Documento} &
Unidad de archivo cargada al sistema, con metadatos de origen, fecha y categoría. \\

\textbf{Fragmento} &
Porción de un documento resultante de la segmentación, identificada de forma única y
localizable en su documento de origen. Es la unidad mínima de evidencia. \\

\textbf{Curación} &
Proceso de seleccionar, comparar, depurar, actualizar y organizar el corpus. Esta
investigación aborda su fase \emph{detectable} (identificar relaciones problemáticas)
y deja fuera su fase \emph{deliberativa} (decidir qué se conserva). \\

\textbf{Defecto} &
Relación problemática entre dos o más fragmentos del corpus, de una de las clases
definidas en §\ref{sec:conc-defectos}. La unidad no es el fragmento sino la relación. \\

\textbf{Sugerencia} &
Propuesta generada por el sistema sobre un defecto detectado, en estado pendiente,
acompañada de los identificadores de los fragmentos que la sustentan. \\

\textbf{Decisión} &
Acto del docente sobre una sugerencia: aceptar, modificar o rechazar. Queda
registrada con su autor, su momento y, en caso de rechazo, su motivo. \\

\textbf{Publicación} &
Incorporación efectiva de un cambio al corpus. Por restricción arquitectónica,
requiere una decisión de aceptación previa. \\
\end{xltabular}

\subsection{Clases de defecto}
\label{sec:conc-defectos}

La taxonomía se deriva de la literatura y no de la implementación: cada clase remite
a un fundamento y admite una definición que permite anotar el \emph{ground truth}.

\begin{xltabular}{\textwidth}{@{}L{3.2cm} X L{3cm}@{}}
\toprule
\textbf{Clase} & \textbf{Definición operacional} & \textbf{Fundamento} \\
\midrule
\endfirsthead
\toprule
\textbf{Clase} & \textbf{Definición operacional} & \textbf{Fundamento} \\
\midrule
\endhead
\bottomrule
\endfoot

\textbf{Contradicción} &
Dos fragmentos del corpus afirman proposiciones incompatibles: un valor numérico, una
definición, un procedimiento, un plazo o un criterio. &
\textcite{koreeda2021contractnli}, \textcite{thorne2018fever} \\

\textbf{Redundancia} &
Dos fragmentos expresan el mismo contenido con redacciones distintas, sin que ninguno
contradiga al otro. &
\textcite{gao2024retrieval}, \textcite{reimers2019sentencebert} \\

\textbf{Desactualización} &
Un fragmento afirma algo que otro más reciente sustituye, o contiene información
específica de una instancia pasada (fechas, cupos, nombres). &
\textcite{leacock2007framework}, \textcite{gao2024retrieval} \\

\textbf{Sin evidencia} &
No existe en el corpus material suficiente para decidir el estatus de una afirmación.
Es un resultado, no una ausencia de resultado. &
\textcite{thorne2018fever}, \textcite{gubelmann2024capturing} \\

\textbf{Control negativo} &
Par de fragmentos temáticamente próximos y no problemáticos, incluido
deliberadamente para estimar la tasa de falsos positivos. &
\textcite{gubelmann2024capturing} \\
\end{xltabular}

\begin{advertencia}
La inclusión de controles negativos no es opcional.
\textcite{gubelmann2024capturing} muestran que, con un desbalance realista donde la
mayoría de los pares no guarda relación, un modelo puede colapsar al atribuir
relaciones inexistentes. Un benchmark compuesto solo por defectos plantados mide
\emph{recall} y no dice nada sobre precisión en condiciones de uso.
\end{advertencia}

\subsection{Términos de evaluación}
\label{sec:conc-evaluacion}

\begin{xltabular}{\textwidth}{@{}L{3.6cm} X@{}}
\toprule
\textbf{Término} & \textbf{Definición operacional} \\
\midrule
\endfirsthead
\toprule
\textbf{Término} & \textbf{Definición operacional} \\
\midrule
\endhead
\bottomrule
\endfoot

\textbf{Verdad de referencia} &
Etiqueta asignada por anotadores humanos a cada instancia del benchmark, con la
evidencia esperada asociada y el acuerdo entre anotadores reportado. \\

\textbf{Detección correcta} &
Coincidencia simultánea de la clase asignada y de la evidencia citada con las de la
verdad de referencia, según el criterio de correspondencia definido en la Parte~V.
Acertar la clase con evidencia equivocada no cuenta como detección correcta
\parencite{thorne2018fever}. \\

\textbf{Fidelidad} &
Proporción de afirmaciones de una salida que están sustentadas por la evidencia
citada \parencite{es2024ragas,min2023factscore}. Se evalúa por afirmación atómica, no
por respuesta completa. \\

\textbf{Relevancia del contexto} &
Proporción de los fragmentos recuperados que son pertinentes para la consulta
\parencite{saadfalcon2024ares}. \\

\textbf{Error de calibración} &
Discrepancia agregada entre la confianza estimada y la frecuencia observada de
acierto, medida sobre un conjunto independiente
\parencite{guo2017calibration}. \\

\textbf{Configuración de comparación} &
Variante del sistema con componentes deshabilitados o sustituidos, usada para atribuir
el desempeño a componentes identificables. \\

\textbf{Juez automático} &
Modelo empleado para evaluar salidas. Su validez no se asume: se establece
comparándolo con anotación humana mediante un coeficiente corregido por azar
\parencite{thakur2025judging}. \\
\end{xltabular}

\subsection{Distinciones que el documento mantiene}
\label{sec:conc-distinciones}

\begin{table}[H]
\centering
\small
\caption{Pares de términos que no deben usarse como equivalentes}
\label{tab:distinciones}
\begin{tabularx}{\textwidth}{@{}L{3.4cm} L{3.4cm} X@{}}
\toprule
\textbf{No confundir} & \textbf{Con} & \textbf{Por qué} \\
\midrule
Similitud semántica & Confianza &
La similitud indica proximidad temática entre textos; no informa sobre la corrección
de una decisión \parencite{reimers2019sentencebert}. \\
Confianza estimada & Calibración &
Un puntaje normalizado no es una probabilidad. La calibración se demuestra contra la
tasa de acierto en datos independientes \parencite{guo2017calibration}. \\
Contradicción & Sin evidencia &
Son clases distintas con dificultades distintas y consecuencias distintas para el
docente \parencite{thorne2018fever,gubelmann2024capturing}. \\
Alucinación de fidelidad & Alucinación de factualidad &
La primera es inconsistencia con el contexto provisto; la segunda, con el mundo. Este
proyecto mide la primera \parencite{huang2025survey}. \\
Trazabilidad del proceso & Fundamentación de la salida &
Registrar qué fragmentos se consultaron no demuestra que la salida se apoye en ellos
\parencite{moreau2013provdm}. \\
Aceptación de una sugerencia & Confianza apropiada &
Aceptar puede reflejar sobredependencia. La dependencia se usa con frecuencia como
sustituto de la confianza sin justificarlo \parencite{mehrotra2024systematic}. \\
Recuperación de precedentes & Aprendizaje &
La memoria de retroalimentación opera en inferencia y no modifica parámetros
\parencite{fernandes2023bridging}. \\
Capacidad técnica & Desempeño &
Que el sistema ejecute una operación no informa sobre con qué exactitud la ejecuta
\parencite{hevner2004design}. \\
\bottomrule
\end{tabularx}
\end{table}
